\section{Introduction}

NOTE: THIS IS A VERY ROUGH DRAFT OF THE INTRODUCTION. IT NEEDS TO BE DISCIPLINED AND REFINED, BUT I WANTED TO GET SOME OF THE KEY IDEAS DOWN ON PAPER FIRST.


Drug overdose deaths remain a pressing issue in the US. The reminants of an opioid crisis persist and ever more headlines highlight the proliferation of drugs like fentanyl and rising overdose deaths from products laced with it, so much so that in 2020, drug overdoses were the leading cause of injury-related death in the country, with over 90,000 fatalities \parencite{hedegaardDrugOverdoseDeaths2020a}. Curbing the misappropriation of prescription drugs and the subsequently associated overdose deaths continues to be a major public health priority.

There is a lot that happens before a drug falls into the hands of an individual who may misuse it and ultimately overdose. In the enterprise of designing a policy intervention to prevent drug overdose deaths, it is necessary to consider how a drug may get to those hands to begin with. To that end, there is not much one can do to curb an individual's actions, the illict/illegal creation and distribution of drugs, nor much we can do to steer the innovative research that pharamaceutical firms undertake without potentially severe economic consequences. However, we can influence the regulatory institutions that govern the approval of drugs and the incentives they create for pharmaceutical firms. In other words, it is necessary to consider the source of entry for drugs into the market and how regulatory institutions shape that entry.

The U.S. Food and Drug Administration (FDA) plays a central role in regulating pharmaceutical markets by determining which drugs may enter the U.S. market. Over the past century the agency has developed a reputation as the primary gatekeeper of drug safety and efficacy, exercising significant influence over pharmaceutical innovation and public health outcomes. As emphasized by \textcite{carpenterINTRODUCTIONGatekeeper2010}, the FDA’s authority derives not only from its statutory powers but also from its institutional reputation for scientific rigor and vigilance against unsafe drugs. Beyond its domestic role, the FDA has also become one of the most influential pharmaceutical regulators in the world. Because access to the U.S. pharmaceutical market is economically critical for drug manufacturers, regulatory decisions made by the FDA often shape the global development and approval strategies of pharmaceutical firms. As a result, the agency is frequently regarded as a benchmark for regulatory rigor, and regulatory agencies in other countries often observe, emulate, or coordinate with FDA standards and practices \parencite{carpenterAmericanPharmaceuticalRegulation2010}.


To that end, the FDA has instituted a lot of policy changes to accelerate the approval of drugs. The Prescription Drug User Fee Act (PDUFA) in 1992 is one of the most significant reforms in this regard. The PDUFA allowed the FDA to collect user fees from pharmaceutical manufacturers to fund the drug review process, with the goal of accelerating regulatory review times and bringing new drugs to market more quickly. A large literature has examined the consequences of this policy, focusing primarily on how faster approval times affect drug safety, innovation incentives, and welfare outcomes \parencite{olsonRiskWeBear2008}.




\TODO{Discipline the introduction}{Merge the intro draft above with the one below.}


The U.S. Food and Drug Administration (FDA) plays a central role in regulating pharmaceutical markets by determining which drugs may enter the U.S. market. Over the past century the agency has developed a reputation as the primary gatekeeper of drug safety and efficacy, exercising significant influence over pharmaceutical innovation and public health outcomes. As emphasized by \textcite{carpenterINTRODUCTIONGatekeeper2010}, the FDA’s authority derives not only from its statutory powers but also from its institutional reputation for scientific rigor and vigilance against unsafe drugs. Beyond its domestic role, the FDA has also become one of the most influential pharmaceutical regulators in the world. Because access to the U.S. pharmaceutical market is economically critical for drug manufacturers, regulatory decisions made by the FDA often shape the global development and approval strategies of pharmaceutical firms. As a result, the agency is frequently regarded as a benchmark for regulatory rigor, and regulatory agencies in other countries often observe, emulate, or coordinate with FDA standards and practices \parencite{carpenterAmericanPharmaceuticalRegulation2010}.

Despite this reputation for scientific rigor and regulatory vigilance, the FDA’s regulatory mandate has long generated policy debate. At the center of these debates is a fundamental tension in pharmaceutical regulation: how to balance rapid access to new and affordable treatments with the agency’s responsibility to ensure drug safety and efficacy. Efforts to improve access have historically taken the form of major legislative reforms. For example, the \href{https://www.congress.gov/bill/98th-congress/house-bill/3605}{Drug Price Competition and Patent Term Restoration Act of 1984} (commonly known as the \href{https://www.fda.gov/drugs/abbreviated-new-drug-application-anda/hatch-waxman-letters}{Hatch–Waxman Act}) fundamentally reshaped the U.S. pharmaceutical market by creating the modern \href{https://www.fda.gov/drugs/types-applications/abbreviated-new-drug-application-anda}{abbreviated new drug application} (ANDA) pathway, dramatically expanding the entry of generic drugs while attempting to preserve incentives for brand-name innovation. Although these reforms substantially improved drug affordability and access, scholars have also noted that they may introduce new regulatory challenges and incentive distortions within pharmaceutical markets \parencite{mehlHatchWaxmanActMarket2006}. These recurring policy debates illustrate the complex set of economic incentives surrounding pharmaceutical regulation and highlight the broader challenge facing the FDA: designing regulatory institutions that simultaneously promote innovation, ensure safety, and expand patient access to medicines.

These debates over regulatory speed and safety played a central role in the policy reforms that culminated in the \href{https://www.fda.gov/industry/fda-user-fee-programs/prescription-drug-user-fee-amendments}{Prescription Drug User Fee Act (PDUFA)} in 1992. The PDUFA allowed the FDA to collect user fees from pharmaceutical manufacturers to fund the drug review process, with the goal of accelerating regulatory review times and bringing new drugs to market more quickly. A large literature has examined the consequences of this policy, focusing primarily on how faster approval times affect drug safety, innovation incentives, and welfare outcomes \parencite{olsonRiskWeBear2008}.

At the same time, the public health stakes of regulatory incentives in pharmaceutical approval are particularly salient in markets for controlled substances. Prescription opioids, stimulants, and other scheduled drugs play an important role in clinical care but also carry substantial abuse potential. Over the past two decades, prescription drugs have played a central role in the rise of drug overdose mortality in the United States, frequently exceeding deaths associated with heroin or cocaine in certain periods \parencite{hedegaardDrugOverdoseDeaths2020a}. In response, regulators and pharmaceutical manufacturers have increasingly explored technological and regulatory strategies designed to mitigate abuse risk, including the development of abuse-deterrent formulations intended to make drugs more difficult to manipulate or misuse \parencite{colemanReducingAbusePotential2010}. However, recent empirical work suggests that supply-side interventions in drug markets can generate unintended behavioral responses. For example, reforms that introduced abuse-deterrent formulations of opioid medications reduced certain forms of prescription misuse but also contributed to substitution toward illicit alternatives such as heroin \parencite{alpertSupplySideDrugPolicy2018a}. These findings highlight the broader economic challenge of regulating pharmaceutical supply: policies that alter the availability or characteristics of legal drugs may reshape incentives and access throughout both legal and illicit markets.

However, relatively little attention has been paid to how changes in approval speed may affect the
\textit{composition} of drugs that reach the market. If expedited review mechanisms disproportionately benefit certain types of pharmaceutical products—particularly those with high expected market returns—then regulatory reforms designed to accelerate approval could alter not only the number of drugs approved, but also the types of drugs entering the market. In particular, pharmaceutical firms producing highly profitable drug classes may have stronger incentives to utilize expedited regulatory pathways.

This paper investigates whether policies designed to accelerate FDA drug approvals have changed the composition of approved drugs in the United States. Specifically, it examines whether the introduction of PDUFA is associated with an increase in the approval of controlled substances relative to other pharmaceutical products. Because controlled substances—such as opioid analgesics—often generate substantial market revenues while also posing significant public health risks, understanding how regulatory institutions shape their entry into the market is of particular policy relevance.

More broadly, examining how regulatory reforms affect the composition of drug approvals can shed light on the economic incentives embedded within pharmaceutical regulation. If policies intended to improve regulatory efficiency systematically favor certain categories of drugs, then regulatory institutions may influence not only the speed of innovation, but also the direction of pharmaceutical development.
